\subsectionaddtolist{الف}
\[
L = \{x \in \{a,b\}^* \mid n_b(x) = n_a(x)^2\}
\]

 این زبان مستقل از متن {نیست}. رابطه بین تعداد حروف $a$ و $b$ به صورت مجذوری (درجه دو) است که فراتر از توانایی حافظه پشته‌ای (PDA) می‌باشد.

فرض می‌کنیم $L$ یک زبان مستقل از متن باشد و $p$ طول پامپینگ آن باشد. 
رشتهٔ $s = a^p b^{p^2}$ را انتخاب می‌کنیم. مشخص است که $s \in L$ زیرا $n_a(s) = p$ و $n_b(s) = p^2$ است و طول رشته $|s| = p + p^2 \ge p$ می‌باشد.

طبق لم پامپینگ، رشته $s$ را می‌توان به صورت $s = uvxyz$ تجزیه کرد به طوری که:
\begin{enumerate}
    \item $|vxy| \le p$
    \item $|vy| \ge 1$
    \item به ازای هر $i \ge 0$، رشته $u v^i x y^i z \in L$ باشد.
\end{enumerate}

با توجه به اینکه $|vxy| \le p$، بخش $vxy$ نمی‌تواند همزمان شامل هر دو بخش ابتدایی $a$ها و بخش انتهایی $b$ها به میزان زیاد باشد تا توازن درجه دو حفظ شود. حالت‌های زیر را بررسی می‌کنیم:
\begin{itemize}
    \item \textbf{حالت ۱: $vy$ فقط شامل حرف $a$ باشد.}
    اگر $i=2$ قرار دهیم، تعداد $a$ها افزایش می‌یابد ($n_a > p$) در حالی که تعداد $b$ها همان $p^2$ باقی می‌ماند. از آنجا که $(p+k)^2 > p^2$ (برای $k \ge 1$)، رابطه $n_b = n_a^2$ برقرار نخواهد بود.
    
    \item \textbf{حالت ۲: $vy$ فقط شامل حرف $b$ باشد.}
    با تغییر $i$، تعداد $b$ها تغییر می‌کند اما تعداد $a$ها همان $p$ باقی می‌مانم، پس تعداد $b$ها دیگر برابر با $p^2$ نخواهد بود.
    
    \item \textbf{حالت ۳: $vy$ شامل هر دو حرف $a$ و $b$ باشد.}
    از آنجا که حروف $a$ قبل از $b$ آمده‌اند، $v$ باید شامل $a$ و $y$ باید شامل $b$ باشد. فرض کنید $v = a^k$ و $y = b^m$ که $k, m \ge 1$ و $k+m \le p$.
    اگر $i=2$ را پمپ کنیم، رشته جدید دارای $p+k$ تعداد $a$ و $p^2+m$ تعداد $b$ خواهد بود. برای اینکه رشته جدید در زبان باشد باید داشته باشیم:
    \[ p^2 + m = (p+k)^2 = p^2 + 2pk + k^2 \implies m = 2pk + k^2 \]
    اما می‌دانیم $m \le p$ و از طرفی چون $k \ge 1$، مقدار $2pk + k^2 \ge 2p + 1$ است. بنابراین $m$ نمی‌تواند برابر با $2pk + k^2$ باشد (تناقض).
\end{itemize}
بنابراین فرض خلف باطل است و زبان $L$ مستقل از متن نیست.

\subsectionaddtolist{ب}
\[
L = \{x \in \{a,b,c\}^* \mid n_a(x) = \min\{n_b(x), n_c(x)\}\}
\]

می‌دانیم اگر $L$ مستقل از متن باشد، اشتراک آن با هر زبان منظم مانند $L_R = a^* b^* c^*$ نیز باید مستقل از متن باشد. نام این زبان مشترک را $L'$ می‌گذاریم:
\[ L' = L \cap a^* b^* c^* = \{a^i b^j c^k \mid i = \min(j, k)\} \]

حالا عدم مستقل از متن بودن $L'$ را با لم پامپینگ اثبات می‌کنیم. فرض کنید $p$ طول پامپینگ باشد. رشتهٔ $s = a^p b^p c^p$ را انتخاب می‌کنیم. در این رشته $n_a = p$ و $\min(n_b, n_c) = \min(p, p) = p$ است، پس $s \in L'$ و $|s| = 3p \ge p$.

رشته $s = uvxyz$ را به صورتی تجزیه می‌کنیم که $|vxy| \le p$ و $|vy| \ge 1$. به دلیل محدودیت طول $|vxy| \le p$، رشته $vy$ حداکثر می‌تواند شامل دو نوع حرف متوالی باشد (یعنی یا شامل $a,b$ است یا شامل $b,c$ و یا فقط شامل یکی از حروف).
\begin{itemize}
    \item \textbf{اگر $vy$ شامل حرف $c$ نباشد} (یعنی فقط شامل $a$ و $b$ باشد):
    با پمپ کردن به ازای $i=2$، تعداد $a$ها یا $b$ها (یا هر دو) افزایش می‌یابد در حالی که تعداد $c$ها همان $p$ باقی می‌ماند. اگر تعداد $a$ها زیاد شود ($n_a > p$)، دیگر برابر با $\min(n_b, n_c) \le p$ نخواهد بود.
    
    \item \textbf{اگر $vy$ شامل حرف $a$ نباشد} (یعنی فقط شامل $b$ و $c$ باشد):
    با پمپ کردن به ازای $i=0$ (حذف کردن)، تعداد $a$ها برابر با $p$ باقی می‌ماند اما تعداد $b$ها یا $c$ها کاهش می‌یابد. بنابراین حداقل یکی از مقادیر $n_b$ یا $n_c$ کمتر از $p$ می‌شود، پس $\min(n_b, n_c) < p$ خواهد شد که با $n_a=p$ برابر نیست.
\end{itemize}
در تمامی حالت‌ها به تناقض می‌رسیم. بنابراین $L'$ و در نتیجه زبان اصلی $L$ مستقل از متن نیستند.

\subsectionaddtolist{ج}
\[
L = \{w \in \{a,b\}^* \mid 2n_a(w) \le n_b(w) \le 3n_a(w)\}
\]

 این زبان {مستقل از متن است}. ساختار این زبان به ما اجازه می‌دهد که به ازای هر حرف $a$، بین ۲ تا ۳ حرف $b$ در هر جایگاهی از رشته داشته باشیم. برای اثبات، گرامر مستقل از متن آن را طراحی می‌کنیم.

قوانین تولید گرامر به صورت زیر تعریف می‌شوند:
\[
S \rightarrow SS \mid aSbSbS \mid aSbSbSbS \mid bSaSbS \mid bSbSaSbS \mid bSbSbSaS \mid \epsilon
\]

به طور شهودی، هر بار که یک $a$ تولید می‌شود، پشته می‌تواند تعادل خود را با دریافت ۲ یا ۳ عدد $b$ حفظ کند. چون ترتیب حروف مهم نیست، گرامر بالا با ترکیب نمادها ساختار مستقل از متن بودن زبان را اثبات می‌کند.
